\documentclass{article}
\PassOptionsToPackage{numbers, sort, compress}{natbib}
\usepackage[main,final]{neurips_2025}

\usepackage[T1]{fontenc}
\usepackage[utf8]{inputenc}
\usepackage{hyperref}
\usepackage{url}
\usepackage{booktabs}
\usepackage{nicefrac}
\usepackage{microtype}
\usepackage{xcolor}
\usepackage{subcaption}
\usepackage{graphicx}
\usepackage{multirow}
\usepackage{amsmath,amssymb,amsfonts}
\usepackage{bbm}
\usepackage{amsthm}

\newtheorem{theorem}{Theorem}
\newtheorem{proposition}[theorem]{Proposition}
\newtheorem{lemma}{Lemma}
\newtheorem{definition}{Definition}
\newtheorem{remark}{Remark}

\title{Image Plagiarism Detection with Equivariant Encoders}

\author{%
  Denis Gudkov\\
  Intelligent Systems Department, MIPT\\
  \texttt{gudkov.dv@phystech.edu}
}

\begin{document}
\maketitle

\begin{abstract}
A pirated image is rarely a pixel-perfect copy: rotations, reflections, crops, and re-compression are routine, and a detector should map all of them to the same point in feature space. We ask whether building these symmetries into the encoder of a siamese plagiarism detector is preferable to approximating them with data augmentation. Three encoders are compared inside one pipeline: a non-equivariant ViT-Tiny/16 baseline trained with geometric augmentation, an Octic ViT with exact dihedral invariance under multiples of ninety-degree rotation and reflection, and a Harmformer with continuous roto-translational invariance covering arbitrary rotation angles and planar translations. Training pairs are formed on COCO\,2017 with a cross-batch protocol; evaluation uses 1\,200 cross-domain pairs sampled from DomainNet across six visually distinct domains. We report false positive rate and recall stratified by transformation class. Results are pending; the experimental protocol is fixed.
\end{abstract}

\textbf{Keywords:} image plagiarism detection; siamese networks; equivariant Vision Transformers; group invariance; contrastive learning; DomainNet.

\section{Introduction}\label{sec:intro}

Two images depict the same content under a plagiarist's edits when one differs from the other by a rotation, a reflection, a crop, or a photometric perturbation, and a detector should map both to the same point in feature space. The question we study is whether to enforce this invariance through the encoder architecture or to approximate it through geometric augmentation at training time. The detector receives a pair of RGB images and returns a similarity score in the unit interval. Positive pairs are produced by composing a geometric transformation, drawn either from the dihedral group of the square or from the special Euclidean group of the plane, with a photometric perturbation that covers brightness and contrast jitter, blur, additive noise, JPEG re-compression, and watermark overlays. Negative pairs consist of two images sampled independently from the same visual domain. The encoder is shared between the two siamese branches and the head is symmetric in its arguments, so that any geometric variant of an image arrives at the same comparison.

Standard Vision Transformers~\citep{dosovitskiy2021image} provide no geometric guarantee. Patch tokenization commutes with neither rotation nor reflection, and learned positional encodings do not extrapolate beyond the training distribution; geometric augmentation can only cover the orbit of a transformation group stochastically, inflates training cost, and makes no promise on transformations unseen during training. Scale alone does not close this gap, since spatial reasoning of large vision models has been shown to fail to improve with parameter count~\citep{chen2024spatialvlm}. Geometric deep learning~\citep{bronstein2021geometric, cohen2016group, weiler2019general} formalizes the alternative, in which the desired invariance is built into the architecture as commutation with a group action and then propagates exactly through the network. Two recent ViT variants make this concrete on transformer backbones. The Octic ViT~\citep{nordstrom2025octic} enforces exact dihedral equivariance through group-structured weight sharing on the eight orientations of the square and reduces FLOPs by forty percent relative to a standard ViT of comparable expressivity. Harmformer~\citep{karella2024harmformer} replaces patch projections with circular harmonic filters so that self-attention commutes with continuous in-plane rotations and translations, and therefore covers all rotation angles, not only multiples of ninety degrees. The siamese template, on which the detector is built, is due to~\citet{bromley1993signature}; the closest plagiarism-specific baseline is the pairwise ViT detector of~\citet{dorin2024pairwise}, which fixes our training protocol but offers no invariance guarantee.

The closest alternatives illustrate the trade-off. A non-equivariant siamese ViT, as in~\citet{dorin2024pairwise}, learns geometric robustness only in expectation over the augmentation distribution, so any orbit element under-sampled at training time becomes a candidate failure mode at test time. Classical near-duplicate detection based on local descriptors and geometric verification~\citep{ke2004efficient} is invariant by construction but degrades under domain shift and photometric edits, where transformer-based features are stronger. Among architectural alternatives, the Octic ViT~\citep{nordstrom2025octic} guarantees invariance to rotations by multiples of ninety degrees and to reflections, but covers no other angle, while Harmformer~\citep{karella2024harmformer} extends the guarantee to all rotation angles at a higher compute cost. None of these architectures has been evaluated on pairwise plagiarism detection, which is the specific gap we address.

The contribution is a single siamese pipeline that wraps three encoders, a non-equivariant ViT-Tiny/16 baseline, an Octic ViT, and a Harmformer, and pits architectural group invariance against geometric augmentation under matched conditions. Two elements distinguish the setup from~\citet{dorin2024pairwise}. The first is a contrastive regularizer based on the squared Euclidean distance between embeddings~\citep{hadsell2006dimensionality, chopra2005learning}, which depends only on Euclidean distances in the feature space and is therefore preserved under any orthogonal change of basis on that space; the property aligns the loss geometry with the geometry of group-invariant encoders. The second is a controlled augmentation ablation: each equivariant encoder is trained without group augmentation, while the ViT-Tiny/16 baseline is trained with the corresponding group-specific augmentation block, so that any difference in error attributes to architectural inductive bias rather than to data exposure. The computational experiment estimates whether group-invariant encoders trained without the group augmentation block match or exceed the augmented baseline on false positive rate and per-class recall on a balanced cross-domain test of 1\,200 pairs sampled from DomainNet~\citep{peng2019domainnet}; training pairs are constructed on COCO\,2017 with the cross-batch protocol of~\citet{dorin2024pairwise}.

The remainder of the paper proceeds as follows. Section~\ref{sec:rw} reviews siamese detectors, equivariant networks, and equivariant ViTs. Section~\ref{sec:problem} formalizes the detection task and states the equivariance hypothesis. Section~\ref{sec:prelim} fixes notation for group actions, representations, and equivariance. Section~\ref{sec:method} describes the three encoders, the symmetric fusion head, and the loss. Section~\ref{sec:exp} reports the experimental protocol with placeholder results. Section~\ref{sec:concl} concludes.

\section{Related Work}\label{sec:rw}

Twin networks with shared weights and a learned distance were introduced by~\citet{bromley1993signature} for signature verification and have since been adopted across pairwise matching problems. Classical near-duplicate detection combined local descriptors with geometric verification~\citep{ke2004efficient}, which gave geometric robustness by construction but was brittle under domain shift and photometric edits, where modern visual features are stronger. \citet{dorin2024pairwise} applied a siamese ViT with a contrastive loss to image plagiarism on DomainNet and proposed the cross-batch pairing protocol that we reuse without modification; the pipeline reports FPR as the primary metric, but offers no invariance guarantee against geometric perturbations and depends on geometric augmentation for robustness.

Group-equivariant convolutional networks~\citep{cohen2016group} encode discrete symmetries through weight sharing on group orbits, and \citet{weiler2019general} extend the construction to the continuous Euclidean group $E(2)$ via steerable filters. \citet{bronstein2021geometric} treat equivariance as an inductive bias that improves sample efficiency under group-structured data, and the same line of reasoning motivates the transformer variants used here. The Octic ViT~\citep{nordstrom2025octic} replaces standard transformer layers with $D_4$-equivariant ones; by Schur's lemma, equivariant linear maps decompose blockwise across the five irreducible representations of $D_4$, which both constrains parameters and reduces FLOPs by 40\% at expressivity comparable to a non-equivariant ViT, but the resulting invariance is restricted to multiples of $90^\circ$ and reflections. Harmformer~\citep{karella2024harmformer} replaces patch projections with circular harmonic filters and combines them with self-attention to obtain continuous roto-translational equivariance over $SE(2)$, which covers every rotation angle but at a higher compute cost than $D_4$. Neither architecture has been evaluated on pairwise plagiarism detection, which leaves open whether the geometric inductive bias they provide actually reduces false positives on this specific task.

The contrastive component of our loss follows the formulation of~\citet{hadsell2006dimensionality} and the metric-learning template of~\citet{chopra2005learning}: it depends only on Euclidean distances between embeddings, which means that an orthogonal change of basis on the embedding space leaves it invariant, and that property aligns the loss geometry with the geometry of $G$-invariant encoders.

\section{Problem Statement}\label{sec:problem}

We treat image plagiarism detection as a binary classification problem on image pairs: given a pair $(x_1, x_2)\in\mathcal{X}\times\mathcal{X}$, predict $y \in \{0,1\}$, where $y=1$ marks $x_2$ as a plagiarized derivative of $x_1$. Plagiarism is operationalized as the composition of a geometric transformation $g \in G$ with a photometric perturbation $\eta$, producing $x_2 = (g \circ \eta)(x_1)$. The geometric component is drawn from one of the symmetry groups $G \in \{D_4, SE(2)\}$, where $D_4$ contains the rotations of the square by multiples of $90^\circ$ and the four reflections, and $SE(2)$ contains all in-plane rotations and translations. The photometric component covers brightness and contrast jitter, Gaussian blur, additive Gaussian noise, JPEG re-compression, and watermark overlays, and is shared across all training and evaluation regimes.

The labelled dataset is $\mathcal{D} = \{(x_{1,i}, x_{2,i}, y_i)\}_{i=1}^{N}$ with $x_{\cdot,i} \in \mathcal{X} = [0,1]^{3\times 224\times 224}$ and $y_i \in \{0,1\}$, balanced between positive and negative pairs and stratified by visual domain. Positive pairs are drawn from the joint distribution $p_+(x_1, x_2) = p(x_1)\,p_T(x_2 \mid x_1)$, where $p_T(x_2 \mid x_1) = \int \delta\bigl(x_2 - (g\circ\eta)(x_1)\bigr)\,d\mu(g, \eta)$ and $\mu$ is a fixed mixture over $G \cup G_{\mathrm{photo}}$. Negative pairs are independent draws from the marginal of the same domain, $p_-(x_1, x_2) = p(x_1)\,p(x_2)$, which produces hard negatives with shared visual statistics rather than easy contrasts across unrelated images. All images are resized to $224\times 224$ and normalised per channel using ImageNet statistics; the test pool spans six DomainNet~\citep{peng2019domainnet} domains: real, painting, clipart, quickdraw, infograph, and sketch.

The detector is the composition $f_\theta(x_1, x_2) = \sigma\bigl(h_\psi(\phi_\theta(x_1), \phi_\theta(x_2))\bigr)$, where $\phi_\theta\colon\mathcal{X}\to\mathbb{R}^d$ is a shared encoder, $h_\psi$ is a fusion head, and $\sigma$ is the sigmoid. Two structural restrictions are placed on the model class. First, the encoder is $G$-invariant, so that $\phi_\theta(g\cdot x) = \phi_\theta(x)$ for all $g\in G$, which implies that all geometric variants of an image collapse to one embedding without augmentation. Second, the head is symmetric, $h_\psi(z_1, z_2) = h_\psi(z_2, z_1)$, so that the score does not depend on the order in which the two images are presented. Training freezes the encoder backbone and updates only the fusion head, which limits the source of variance to the head and to augmentation sampling and reflects a realistic compute budget on a single GPU with at most 48\,GB of VRAM.

The loss combines class-weighted binary cross-entropy with an $L^2$ contrastive regularizer,
\begin{equation}\label{eq:loss}
  \mathcal{L} \;=\; \mathcal{L}_{\mathrm{BCE}}(w_+, w_-) \;+\; \lambda\,\mathcal{L}_{\mathrm{contr}}^{L^2},
  \qquad
  \mathcal{L}_{\mathrm{contr}}^{L^2}(z_1, z_2, y)
    = y\,\|z_1-z_2\|_2^2 + (1-y)\,\max\bigl(0, m - \|z_1-z_2\|_2\bigr)^2,
\end{equation}
with $z_i = \phi_\theta(x_i)$, contrastive margin $m > 0$, and BCE class weights $w_+, w_-$ that compensate for label imbalance inside a cross-batch pairing. Generalisation is estimated by stratified five-fold cross-validation over the DomainNet pool, with stratification by domain so that every fold contains all six domains; the decision threshold is selected on the validation fold by maximizing F1 and applied unchanged on the test fold. The external quality criteria are the false positive rate at the F1-maximizing threshold and Recall stratified by transformation class, both reported as mean and standard deviation across folds; the secondary criterion is F1 together with a t-SNE visualisation of embeddings on the full evaluation pool.

The research hypothesis closes the statement. Let $\phi^G$ denote a $G$-invariant ViT trained without block-$G$ augmentation, and $\phi^{\mathrm{aug}}$ a non-equivariant ViT-Tiny/16 trained with block-$G$ augmentation. For every transformation class $T\subseteq G$ we test
\begin{equation}\label{eq:hypothesis}
  H_1: \quad
  \mathrm{FPR}_T(\phi^G)\le \mathrm{FPR}_T(\phi^{\mathrm{aug}})
  \quad\text{and}\quad
  \mathrm{Recall}_T(\phi^G)\ge \mathrm{Recall}_T(\phi^{\mathrm{aug}}),
\end{equation}
with at least one inequality strict on some $T$. The null hypothesis $H_0$ is the complementary statement that the differences in $\mathrm{FPR}_T$ and $\mathrm{Recall}_T$ do not exceed the cross-fold variability of $\phi^{\mathrm{aug}}$. Optimal parameters are obtained by
\begin{equation}\label{eq:argmin}
  \hat\theta = \arg\min_\theta\ \mathbb{E}_{(x_1, x_2, y)\sim\mathcal{D}}\,\mathcal{L}(\theta),
\end{equation}
approximated by mini-batch stochastic optimization, and the returned checkpoint is the one with lowest validation loss across the five cross-validation folds.

\section{Preliminaries: Equivariance}\label{sec:prelim}

A left action of a group $G$ on a set $\mathcal{X}$ is a map $G\times\mathcal{X}\to\mathcal{X}$, $(g, x)\mapsto g\cdot x$, satisfying $e\cdot x = x$ and $(gh)\cdot x = g\cdot(h\cdot x)$ for all $g, h\in G$ and all $x\in\mathcal{X}$. A linear representation of $G$ on a vector space $V$ is a homomorphism $\rho\colon G\to\mathrm{GL}(V)$, so that group multiplication is preserved by $\rho(gh) = \rho(g)\rho(h)$. A map $\phi\colon\mathcal{X}\to\mathbb{R}^d$ is $G$-equivariant with respect to representations $\rho_{\mathcal{X}}$ on $\mathcal{X}$ and $\rho$ on $\mathbb{R}^d$ if $\phi(\rho_{\mathcal{X}}(g)\,x) = \rho(g)\,\phi(x)$ for every $g\in G$ and every $x\in\mathcal{X}$, and it is $G$-invariant when $\rho$ is the trivial representation, that is, when $\phi(g\cdot x) = \phi(x)$. For plagiarism detection, $G$-invariance of the encoder $\phi_\theta$ implies that all geometric variants of an image collapse to a single embedding without augmentation, so the head receives identical inputs whether or not a $g\in G$ has been applied; the detection problem reduces to comparing two points whose distance no longer depends on the geometric component of the perturbation.

Schur's lemma states that any $G$-equivariant linear map between irreducible representations is either zero or an isomorphism, and a direct consequence is that equivariant linear maps between fully decomposed representations are block-diagonal across the irreducible decomposition: each block acts within one irreducible component and cannot mix distinct irreducible types. The Octic ViT~\citep{nordstrom2025octic} relies on this property, since each transformer layer operates independently on each irreducible component of $D_4$, so the block-diagonal structure holds by construction and the parameter count drops without loss of expressivity inside each isotypic component.

Two groups appear in the experiment. The dihedral group $D_4 = \langle r, s \mid r^4 = s^2 = e,\ srs = r^{-1}\rangle$ is the symmetry group of the square, of order eight, and acts on images by rotating and reflecting the spatial axes by multiples of $90^\circ$; over the reals it has five irreducible representations, four one-dimensional and one two-dimensional, and any $D_4$-representation decomposes as a direct sum across these five types. The special Euclidean group $SE(2) = SO(2) \ltimes \mathbb{R}^2$ is the group of rigid motions of the plane: an element rotates spatial coordinates by an angle $\alpha$ and translates by a vector in $\mathbb{R}^2$. Because $D_4$ covers only multiples of $90^\circ$ while $SE(2)$ admits arbitrary rotation angles, $SE(2)$-equivariance is strictly more general but also harder to realise exactly on a discrete pixel grid, which is the trade-off encoded in the comparison between Octic ViT and Harmformer.

\section{Method}\label{sec:method}

\subsection{Siamese pipeline}

The detector follows the siamese template of~\citet{bromley1993signature}. Two branches share one encoder $\phi_\theta$ and process the query image $x_1$ and the candidate image $x_2$ in parallel to produce embeddings $z_1 = \phi_\theta(x_1)$ and $z_2 = \phi_\theta(x_2)$ in $\mathbb{R}^d$. A fusion module $h_\psi$ collapses the pair into a scalar, and a sigmoid maps it to a plagiarism probability. Weight sharing places $z_1$ and $z_2$ in the same coordinate system, which is the precondition for any pairwise comparison: without it, the two branches would produce vectors in unrelated coordinate systems and the head could only learn an arbitrary alignment.

\subsection{Encoders}

The non-equivariant reference is a ViT-Tiny/16~\citep{dosovitskiy2021image}. A $224\times 224$ input is split into $14\times 14 = 196$ non-overlapping $16\times 16$ patches, each linearly projected to $\mathbb{R}^{192}$, and the resulting sequence with a prepended \texttt{[CLS]} token is processed by twelve transformer blocks with multi-head self-attention. The \texttt{[CLS]} embedding after the final layer norm is taken as the image representation $z\in\mathbb{R}^{192}$. The encoder has no equivariance guarantee, and at training time it receives both the photometric block and a group-specific geometric block matched to the comparison in which it participates.

The Octic ViT~\citep{nordstrom2025octic} replaces each transformer layer with a $D_4$-equivariant variant. By Schur's lemma the embedding space decomposes across the irreducibles of $D_4$,
\begin{equation}\label{eq:octic-decomp}
  \mathbb{R}^d \;\cong\;
  (\mathbb{R}^{d/8})_{A_1}\oplus(\mathbb{R}^{d/8})_{A_2}\oplus(\mathbb{R}^{d/8})_{B_1}\oplus(\mathbb{R}^{d/8})_{B_2}
  \;\oplus\;(\mathbb{R}^{d/4})_{E},
\end{equation}
into four one-dimensional isotypic components and one two-dimensional component. $D_4$-invariant pooling projects onto the trivial isotypic component $A_1$ to produce a $D_4$-invariant output, and the encoder is trained with the photometric block only.

Harmformer~\citep{karella2024harmformer} replaces patch embeddings with circular harmonic filters that act as steerable feature maps and combines them with self-attention to obtain equivariance under continuous roto-translations. Where the Octic ViT is restricted to multiples of $90^\circ$, Harmformer covers arbitrary rotation angles in $SE(2)$, and the stronger guarantee comes at a higher computational cost per layer. Like the Octic ViT, Harmformer is trained with the photometric block only.

\subsection{Fusion head}

The head takes $z_1, z_2 \in \mathbb{R}^d$ and returns a logit through pair features and a two-layer MLP. Set $\delta = |z_1 - z_2|$ for the element-wise absolute difference, $p = z_1\odot z_2$ for the element-wise product, and $v = [u_\delta(\delta);\, u_p(p)]\in\mathbb{R}^{2d}$ for the concatenation after independent linear maps; the head value is then $h_\psi(z_1, z_2) = c^\top \mathrm{ReLU}(A v + b)$. Since $|z_1 - z_2| = |z_2 - z_1|$ and $z_1\odot z_2 = z_2\odot z_1$, the vector $v$ is unchanged under transposition of the inputs, hence the head is symmetric in its two arguments.

\begin{proposition}[Symmetry of the head]\label{prop:head-sym}
$h_\psi(z_1, z_2) = h_\psi(z_2, z_1)$ for all $z_1, z_2\in\mathbb{R}^d$.
\end{proposition}

\subsection{$L^2$ contrastive regularizer}

We use the regularizer in~\eqref{eq:loss}, with $L^2$ distance rather than cosine similarity. Working with $L^2$ distance is deliberate: its value is preserved by any orthogonal change of basis on $\mathbb{R}^d$, since $\|U z_1 - U z_2\|_2 = \|U(z_1 - z_2)\|_2 = \|z_1 - z_2\|_2$ for every orthogonal $U$, while cosine similarity depends on individual norms, which an isometry need not preserve. The practical consequence is that two $G$-invariant encoders that differ only by an orthogonal change of basis on the embedding space yield identical loss values on the same pair, so the regularizer does not penalize one rotation of the embedding frame over another.

\begin{proposition}[Isometry invariance of the regularizer]\label{prop:iso-loss}
For every orthogonal $U\colon\mathbb{R}^d\to\mathbb{R}^d$ and every $(z_1, z_2, y)$,
$\mathcal{L}_{\mathrm{contr}}^{L^2}(U z_1, U z_2, y) = \mathcal{L}_{\mathrm{contr}}^{L^2}(z_1, z_2, y)$.
\end{proposition}

\subsection{Augmentation blocks}\label{sec:aug}

The photometric block is shared across all encoders and all runs and is sampled independently for each pair. It applies, with probability $p = 0.3$ each, ColorJitter on brightness, contrast, saturation, and hue; Gaussian Blur with a small kernel; additive Gaussian Noise; JPEG re-compression at random quality; and a watermark overlay that mimics text-and-logo edits common in image piracy. The geometric block is added only to the ViT-Tiny/16 baseline and is matched to the symmetry group of the comparison: in the $D_4$ comparison it samples rotations from $\{90^\circ, 180^\circ, 270^\circ\}$ together with horizontal and vertical flips, and in the $SE(2)$ comparison it adds random crop and isotropic scale on top of the same set of rotations and flips. Equivariant encoders return identical embeddings under any element of the corresponding group, so applying the geometric block to them generates pairs with zero gradient on the head: harmless but redundant. The experimental claim is therefore that the geometric block can be dropped for $\phi^G$ without loss of detection performance.

\subsection{Invariance propagation}

The two structural choices on the encoder and the head combine into a single property of the detector: the score is invariant under any element of $G$ applied jointly to both inputs. If $\phi_\theta$ is $G$-invariant, then $\phi_\theta(g\cdot x_i) = \phi_\theta(x_i)$ for $i = 1, 2$, and since $h_\psi$ and $\sigma$ are deterministic functions of their inputs, $f_\theta(g\cdot x_1, g\cdot x_2) = f_\theta(x_1, x_2)$ for every $g\in G$ and every pair $(x_1, x_2)$. The plagiarism score is therefore the same whether or not $g$ has been applied, which is the invariance required for geometric plagiarism detection.

\begin{proposition}[Joint invariance]\label{prop:joint}
If $\phi_\theta$ is $G$-invariant, then $f_\theta(g\cdot x_1, g\cdot x_2) = f_\theta(x_1, x_2)$ for all $g\in G$ and all $x_1, x_2\in\mathcal{X}$.
\end{proposition}

\section{Computational Experiment}\label{sec:exp}

The goal of the experiment is to test, for each $G\in\{D_4, SE(2)\}$, whether a $G$-invariant encoder trained without block-$G$ augmentation matches or exceeds a ViT-Tiny/16 baseline trained with block-$G$ augmentation on FPR and per-class Recall. Training pairs are constructed on COCO\,2017 with the cross-batch protocol of~\citet{dorin2024pairwise}: a batch of $B$ images is augmented to two views per image, applying the photometric block to all encoders and adding the geometric block on top for the augmented ViT-Tiny/16 baseline, all $B^2$ ordered pairs are formed with labels $y_{ij} = \mathbbm{1}[i = j]$, embeddings are computed through the frozen encoder and the fusion module, the composite loss in~\eqref{eq:loss} is evaluated, and AdamW updates the fusion parameters under a multi-step learning-rate schedule with class-weighted binary cross-entropy and a fixed contrastive margin.

Evaluation uses the DomainNet benchmark~\citep{peng2019domainnet}, which spans six visually diverse domains: real, painting, clipart, quickdraw, infograph, and sketch. We construct a balanced test set of 1\,200 pairs as 100 positive and 100 negative pairs per domain. A positive pair is a source image together with $(g\circ\eta)(x)$, with $g$ drawn uniformly from the symmetry group of the comparison and $\eta$ from the photometric mixture; a negative pair is two distinct images drawn from the same domain, which produces hard negatives with shared visual statistics. The protocol enables stratified evaluation by transformation class and by domain.

The experiment comprises four training runs grouped into two per-group comparisons. In the first comparison, $G = D_4$, a ViT-Tiny/16 trained with the photometric block plus rotations from $\{90^\circ, 180^\circ, 270^\circ\}$ and the two flips is compared against the Octic ViT trained with the photometric block only. In the second comparison, $G = SE(2)$, a ViT-Tiny/16 trained with the photometric block plus the same rotations and flips together with random crop and isotropic scale is compared against Harmformer trained with the photometric block only. The schema is summarized in Table~\ref{tab:runs-schema}; metric cells are filled after the runs complete.

\begin{table}[h]
\centering
\small
\caption{Run schema. Group augmentations are applied only to ViT-Tiny/16; equivariant encoders see the photometric block only.}
\label{tab:runs-schema}
\begin{tabular}{@{}llcc@{}}
\toprule
Encoder & Group augmentations on ViT-Tiny/16 & FPR & Recall \\
\midrule
\multicolumn{4}{@{}l}{\textit{Comparison 1: $G = D_4$}} \\
\quad ViT-Tiny/16 + block~G$_{D_4}$ & R90/180/270, HFlip, VFlip & --- & --- \\
\quad Octic ViT (no group augs)     & ---                       & --- & --- \\
\midrule
\multicolumn{4}{@{}l}{\textit{Comparison 2: $G = SE(2)$}} \\
\quad ViT-Tiny/16 + block~G$_{SE(2)}$ & R90/180/270, HFlip, VFlip, Crop, Scale & --- & --- \\
\quad Harmformer (no group augs)      & ---                                    & --- & --- \\
\bottomrule
\end{tabular}
\end{table}

For each run the fusion head is trained on the training split, the decision threshold is selected on the validation fold by maximizing F1, and final metrics are computed on the held-out test set. The primary metrics are FPR at the F1-maximizing threshold and Recall, both stratified by transformation class, with full ROC and PR curves and AUC reported per comparison. Stratified Recall directly probes the equivariance hypothesis: a $G$-invariant encoder is expected to be approximately flat across orbit elements of $G$, while an augmentation-only baseline can dip between orbit points whose specific transformation was under-sampled at training time. Code, configuration files, and DomainNet pair lists are released at \url{https://github.com/intsystems/2026-Project-198}.

Results are pending the training runs; the report will be filled in the camera-ready version. The reserved figures are a bar plot of FPR per encoder and augmentation regime, with grouped bars contrasting the configuration with and without the geometric block; a per-transformation Recall heatmap whose rows index encoders and whose columns index transformation classes, including pairwise combinations of the photometric and geometric perturbations; ROC and PR curves per comparison contrasting $\phi^G$ against $\phi^{\mathrm{aug}}$; and a Recall curve as a function of rotation angle $\alpha\in[0^\circ, 360^\circ]$ sampled on a $15^\circ$ grid, on which equivariant encoders are expected to be approximately flat while the augmented ViT-Tiny/16 may dip between $D_4$ orbit points if the augmentation distribution under-samples them. The error report will isolate three sources of variance: stratified errors per transformation class, with differences on R90, R180, R270 directly testing the $D_4$ hypothesis and continuous-angle Recall directly testing the $SE(2)$ hypothesis; cross-fold variance reported as mean and standard deviation over five stratified folds, used to check whether differences between encoders exceed the variability of the baseline; and a failure inspection of the top-$k$ false positives and false negatives per domain, with the applied transformation logged, to attribute residual errors to specific perturbations rather than to encoder choice as a whole.

\section{Conclusion}\label{sec:concl}

We described a siamese plagiarism detector that wraps three encoders inside one pipeline and pits architectural $G$-invariance against geometric augmentation under matched conditions. The setup is fixed: a ViT-Tiny/16 baseline trained with the photometric block plus group-specific augmentations, an Octic ViT for $D_4$, and a Harmformer for $SE(2)$, each evaluated on a balanced cross-domain test of 1\,200 pairs from DomainNet. Three properties of the pipeline follow from the design without invoking the data, namely symmetry of the fusion head, isometry invariance of the $L^2$ contrastive regularizer, and joint invariance propagation through the full siamese score, and they together imply that the plagiarism probability is unchanged whenever a group element is applied jointly to both inputs. Whether $G$-invariant encoders trained without block-$G$ augmentation match or exceed an augmented ViT-Tiny/16 on FPR and Recall is the empirical question left to the DomainNet evaluation. Future work will widen the scope of the geometric component beyond $D_4$ and $SE(2)$ and extend the protocol from frozen backbones to end-to-end training of equivariant ViTs on plagiarism-specific data.

\bibliographystyle{unsrtnat}
\bibliography{references}

\end{document}
